\section{Theoretical Background and Related Work}
\subsection{Shortest-Path Search and Heuristic Assumptions}
For directed $G=(V,E,w)$ with nonnegative weights, Dijkstra selects the smallest tentative start cost~\cite{dijkstra1959}. \astar{} adds domain knowledge~\cite{hart1968}:
\begin{equation}f(n)=g(n)+h(n).
\end{equation}
Here $g(n)$ is the best start-to-$n$ cost found so far and $h(n)$ estimates remaining goal cost. OPEN contains pending states; improvements update predecessors. Acceptance occurs on a valid minimum-priority goal extraction. Setting $h=0$ gives the project's Dijkstra baseline.

Admissibility requires $0\leq h(n)\leq h^*(n)$, with $h^*$ the true remaining graph cost. Consistency requires
\begin{equation}
h(u)\leq w(u,v)+h(v)\quad\text{for every }(u,v)\in E,\qquad h(t)=0.
\label{eq:consistency}
\end{equation}
Consistency makes each valid first expansion optimal. Admissible but inconsistent estimates require reopening improved states. Hart et al.~\cite{hart1968} distinguish these assumptions and tie resolution. Their efficiency analysis is not universal wall-time superiority: heuristic computation and representation also cost time.

\subsection{Geometric Representation and Replanning}
Grids restrict headings and can miss narrow passages. Classical visibility-based planning uses obstacle-free geometric connections~\cite{lozano1979}; finite robot footprints require configuration-space treatment. A complete polygonal visibility graph differs from our sparse scan-center graph. Optimality on the latter does not imply a continuous shortest path.

Theta* examines this tradeoff~\cite{daniel2010}. It propagates grid search information while permitting off-grid headings, but Basic Theta* is not guaranteed to find the true continuous shortest path. Random-grid and game-map experiments compare length, time, expansions, and heading changes: shorter paths than grid smoothing are possible without full visibility-graph construction. Our system instead retains ordinary A* on verified continuous links. The relevant lesson is to assess representation and search separately.

D* Lite reuses previous search information through locally consistent value updates~\cite{koenig2002}. Its unknown-terrain experiments reduce operations relative to repeated search in tested navigation and mapping tasks. Those grids initially treat unknown traversal optimistically; our graph admits only verified free links. Repeated A* is retained for implementation transparency. D* Lite is a relevant alternative, not an unmeasured project speedup.

\subsection{Observation, Exploration, and Escape}
Elfes' occupancy grids accumulate uncertain sensor evidence probabilistically~\cite{elfes1989}. Sonar and stereo demonstrations show the value of additional viewpoints. Our deterministic labels and polygon unions preserve unknown space and accumulation, but omit probabilistic uncertainty and localization drift.

Yamauchi identifies free/unknown frontiers and visits accessible targets~\cite{yamauchi1997}. Office-robot experiments address sensor artifacts; idealized coverage reasoning remains conditional on sensing and reachability. Our polygon candidates similarly seek observations but use ancestry and C++ A* routing. Estimated frontier utility does not prove actual information gain, and one graph-unreachable candidate does not exhaust a branch.

TangentBug combines range discontinuities and a local tangent graph with target motion and boundary following~\cite{kamon1996}. Its convergence analysis assumes static finite-perimeter geometry and ideal point motion; selected simulations compare it with VisBug. Our simpler observed-parent policy does not inherit that guarantee. Together these studies explain why local path optimality requires a separate exploration and escape mechanism.

\subsection{Bundle-Based Planning and the Assigned BAR Problem}
Phan et al.~\cite{phan2026} accumulate sights, rank active open points, construct a skeleton graph, and shorten polylines through ordered bundles of line segments. DAP adjusts vertices on critical segments within represented explored space; Section 3 relates approximation quality to geometric and numerical assumptions. Section 4.5 constructs bundles and Algorithm 2 integrates observations, ranking, refinement, and movement. The main contribution is geometric representation and refinement, not ordinary graph search.

Algorithm 2 is explicitly not generally convergent. Section 5 studies a special algorithm with obstacle-diameter and near-goal connectivity assumptions, which do not prove completion of our maze exploration. Section 6 compares local paths with a grid A* baseline and RRT*, and further evaluates RRTX. Dijkstra is background. The source's \source{Graph.BFS_skeleton_path} nevertheless performs cumulative-cost priority-queue search, despite its BFS-oriented name.

Our evaluation therefore holds observations and endpoints fixed while disclosing different graphs, boundary rules, and motion semantics. The implemented application covers restricted limited-sensing exploration and certified observed-parent return; it reproduces neither bundle optimization nor general polygonal BAR detection. The supplementary literature matrix maps all eight supplied PDFs (seven unique studies) to these themes and manuscript citations, without unsupported novelty claims.
